\subsubsection{Projection of point onto line}

\paragraph{Idea intuitiva.}
Proyeccion ortogonal de $p$ sobre la recta que pasa por $a$ y $b$. El parametro $t$ tal que la proyeccion es $a + t \cdot (b - a)$ es:
\[ t = \frac{(p - a) \cdot (b - a)}{|b - a|^2}. \]
La idea: descomponemos $\vec{ap}$ en una componente paralela a $\vec{ab}$ y una perpendicular; $t$ mide la posicion en la primera.

\paragraph{Propiedades clave (invariantes).}
\begin{itemize}\setlength\itemsep{0pt}
	\item Si $t \in [0, 1]$, la proyeccion cae sobre el \textbf{segmento}; si no, cae sobre la extension.
	\item Retorna \texttt{double} (puede tener precision issues).
	\item $t < 0$: $p$ esta ``antes'' de $a$; $t > 1$: $p$ esta ``despues'' de $b$.
\end{itemize}

\paragraph{Codigo clave.}
Proyeccion ortogonal:
\begin{verbatim}
Point projection(Point p, Point a, Point b) {
    Point ab = b - a, ap = p - a;
    double t = (double)dot(ap, ab) / ab.len2();
    return {a.x + t * ab.x, a.y + t * ab.y};
}
\end{verbatim}

\paragraph{Complejidad.}
\begin{itemize}\setlength\itemsep{0pt}
	\item $O(1)$.
	\item Constante minima (dos dot products y una division).
\end{itemize}

\paragraph{Donde tocar para modificar.}
\begin{itemize}\setlength\itemsep{0pt}
	\item \textbf{Proyeccion sobre segmento}: clipar $t$ a $[0, 1]$ antes de calcular el punto.
	\item \textbf{Distancia al segmento via proyeccion}: combinar con el modulo anterior.
	\item \textbf{Solo el parametro $t$}: si solo necesitas saber si esta dentro del segmento, retorna $t$ sin construir el punto.
\end{itemize}

\paragraph{Trampas y errores frecuentes.}
\begin{itemize}\setlength\itemsep{0pt}
	\item Si $a = b$ ($|ab| = 0$), la proyeccion no esta definida; aniadir guard.
	\item En doubles, aniadir tolerancia \texttt{eps} para comparar $t$ con 0 o 1.
	\item La division por \texttt{len2} puede ser cara si las coordenadas son grandes; usar \texttt{double} o precomputar.
\end{itemize}

\paragraph{Codigo.}
Ver \texttt{cpp.tex}, seccion \textit{Projection of point onto line}.

\vspace{0.5em}
